%---
% title: "Parameter sharing works best across width"
% date: 2026-06-28
% work_order: 6
% permalink: /posts/2026/06/irregular-parameter-sharing-transformers/
% tags:
%   - machine-learning
%   - language-models
%   - transformers
%   - parameter-sharing
%---
\documentclass[11pt]{article}
\usepackage{publication-web}
\begin{document}

Hard parameter sharing reduces model size, but its topology determines which computations are forced to coincide. In a 16-layer, width-1024 Transformer language model, cycle width sharing reached \(4.5043\) validation cross-entropy over three seeds, versus \(4.5645\) for balanced random sharing and \(4.5279\) for maximum-depth-distance sharing. Every shared variant assigned 128 exact MLP blocks to 256 layer-by-width positions, using each block twice. The result suggests preserving layer-specific transformations and sharing parallel width chunks.

\section{The topology question}\label{the-topology-question}

Transformer sharing has usually meant cross-layer reuse: Universal Transformers recur one transition function, ALBERT shares parameters across layers, and later work compares depth schedules (\href{https://arxiv.org/abs/1807.03819}{Dehghani et al.~2018}; \href{https://arxiv.org/abs/1909.11942}{Lan et al.~2019}; \href{https://arxiv.org/abs/2104.06022}{Takase and Kiyono 2021}). The unresolved choice is which computational positions should share under a fixed parameter budget.

Split each of 16 layer MLPs into 16 hidden-width chunks. A schedule assigns each position \((\ell,c)\) to one of \(K=128\) learned blocks,

\[
a(\ell,c)\in\{0,\ldots,127\},
\qquad
\mathrm{MLP}_{\ell}(x)
=\sum_{c=0}^{15}
\phi\!\left(xW^{(1)}_{a(\ell,c)}\right)
W^{(2)}_{a(\ell,c)}.
\]

The tests compare sequence and cycle reuse along depth, sequence and cycle reuse along width, diagonal depth--width reuse, balanced random sharing, maximum-depth-distance sharing, and a best-of-12 random search. The unshared dense model supplies the practical baseline; balanced random sharing is the parameter-matched baseline. The random search selected after 800 proxy steps, then retrained the selected layout for the final three seeds.

\section{Setup and results}\label{setup-and-results}

The decoder-only model used 16 layers, width 1024, 16 heads, sequence length 256, GELU MLPs, and tied output embeddings. Each final run trained for 5,000 AdamW steps on 100M GPT-2-tokenized OpenWebText tokens, with 2M validation tokens and seeds 0--2 (\href{https://cdn.openai.com/better-language-models/language_models_are_unsupervised_multitask_learners.pdf}{Radford et al.~2019}; \href{https://skylion007.github.io/OpenWebTextCorpus/}{Gokaslan and Cohen 2019}). Code and artifacts are in \href{https://github.com/Axym-Labs/irregular-parameter-sharing}{\texttt{Axym-Labs/irregular-parameter-sharing}}.

\publicationfigure[\linewidth]
  {images/parameter-sharing-topology/validation_ce_by_topology.png}
  {/images/parameter-sharing-topology/validation_ce_by_topology.png}
  {Validation cross-entropy by hard-sharing topology}
  {\textbf{Figure 1.} Validation cross-entropy by topology; error bars are across three seeds.}

Table 1. Final validation losses. Shared variants use 128 MLP blocks and 67.1M MLP-bank parameters.

\begin{longtable}[]{@{}
  >{\raggedright\arraybackslash}p{(\linewidth - 10\tabcolsep) * \real{0.1364}}
  >{\raggedright\arraybackslash}p{(\linewidth - 10\tabcolsep) * \real{0.1364}}
  >{\raggedleft\arraybackslash}p{(\linewidth - 10\tabcolsep) * \real{0.1818}}
  >{\raggedleft\arraybackslash}p{(\linewidth - 10\tabcolsep) * \real{0.1818}}
  >{\raggedleft\arraybackslash}p{(\linewidth - 10\tabcolsep) * \real{0.1818}}
  >{\raggedleft\arraybackslash}p{(\linewidth - 10\tabcolsep) * \real{0.1818}}@{}}
\toprule\noalign{}
\begin{minipage}[b]{\linewidth}\raggedright
Variant
\end{minipage} & \begin{minipage}[b]{\linewidth}\raggedright
Sharing topology
\end{minipage} & \begin{minipage}[b]{\linewidth}\raggedleft
MLP blocks
\end{minipage} & \begin{minipage}[b]{\linewidth}\raggedleft
MLP-bank params
\end{minipage} & \begin{minipage}[b]{\linewidth}\raggedleft
Validation CE
\end{minipage} & \begin{minipage}[b]{\linewidth}\raggedleft
Delta vs random
\end{minipage} \\
\midrule\noalign{}
\endhead
\bottomrule\noalign{}
\endlastfoot
Unshared dense MLP & none & 256 & 134.2M & 4.5928 +/- 0.0384 & +0.0284 \\
Sequence depth & depth & 128 & 67.1M & 4.5723 +/- 0.0781 & +0.0078 \\
Cycle depth & depth & 128 & 67.1M & 4.5322 +/- 0.0733 & -0.0322 \\
Sequence width & width & 128 & 67.1M & 4.5621 +/- 0.0366 & -0.0023 \\
Cycle width & width & 128 & 67.1M & \textbf{4.5043 +/- 0.0537} & \textbf{-0.0601} \\
Diagonal depth-width & depth and width & 128 & 67.1M & 4.5504 +/- 0.0086 & -0.0140 \\
Balanced random & depth and width & 128 & 67.1M & 4.5645 +/- 0.0601 & 0.0000 \\
Maximum depth distance & depth and width & 128 & 67.1M & 4.5279 +/- 0.0243 & -0.0366 \\
Best-of-12 random & depth and width & 128 & 67.1M & 4.6057 +/- 0.0523 & +0.0413 \\
\end{longtable}

Cycle width sharing beat balanced random by \(0.0601\) CE and the unshared model by \(0.0885\) within this training budget. Maximum-depth-distance sharing beat random by \(0.0366\). The best proxy among 12 random layouts did not survive final retraining.

\publicationfigure[0.9\linewidth]
  {images/parameter-sharing-topology/delta_vs_random.png}
  {/images/parameter-sharing-topology/delta_vs_random.png}
  {Validation cross-entropy deltas relative to balanced random hard sharing}
  {\textbf{Figure 2.} Parameter-matched topologies relative to balanced random sharing; negative values are better.}

\section{Why width sharing wins}\label{why-width-sharing-wins}

Cycle width sharing keeps a separate block bank in every layer:

\[
a_{\mathrm{width\ cycle}}(\ell,c)=8\ell+(c\bmod 8).
\]

Chunks \(c\) and \(c+8\) share within a layer; no block crosses depth. Its parameter count matches every other shared model.

\publicationfigure[\linewidth]
  {images/parameter-sharing-topology/sharing_topology_schedules.png}
  {/images/parameter-sharing-topology/sharing_topology_schedules.png}
  {Four hard-sharing schedules over layer and MLP chunk positions}
  {\textbf{Figure 3.} Representative schedules over the 16-layer by 16-chunk grid. Color denotes block identity modulo the plotting palette.}

Depth is ordered: lower and upper layers operate on different residual-stream distributions and at different stages of the computation. Width chunks are parallel contributors to one layer update. Cross-width sharing removes capacity while preserving the layer-specific transformation; cross-depth sharing forces one operator into distinct computational regimes.

The depth results support the same distinction. Separating tied uses improves cycle depth over sequence depth, and maximum-depth-distance sharing improves over random. Avoiding cross-depth ties entirely performs best. The supported design rule is therefore about the sharing axis, not the arbitrary ordering of MLP chunks: preserve depth before width.

\end{document}
